% Figura 3: imagem real em um espelho côncavo.
% C na origem, R = 7,84; objeto A a p = 1,5R do vértice, imagem A' a p' = 0,75R
% (1/p + 1/p' = 2/R), altura -0,5 da do objeto. Os raios são desenhados de forma
% idealizada (todos passam por A').
\begin{tikzpicture}[font=\small]
  \def\R{7.84}
  \coordinate (C) at (0,0);
  \coordinate (F) at (\R/2,0);
  \coordinate (A) at (-\R/2,\R/5);
  \coordinate (A') at (\R/4,-\R/10);
  \draw[eixo] (-4.41,0) -- (8.428,0);
  \draw[espelho] (-12.5:\R) arc[start angle=-12.5,end angle=12.5,radius=\R];
  \foreach \a in {-12,-10.5,...,12}
    \draw[hachura] (\a:\R) -- ++(\a-40:0.195);
  \foreach \y/\c in {0.15/azul,0.075/laranja,0/verde,-0.075/rosa,-0.15/vermelhao} {
    \pgfmathsetmacro{\yy}{\y*\R}
    \pgfmathsetmacro{\xx}{sqrt(\R*\R-\yy*\yy)}
    \coordinate (M) at (\xx,\yy);
    \draw[raio,\c] (A) -- (M);
    \draw[raio,\c] (M) -- (A');
    \draw[raio,\c] (A') -- ($(A')!-0.29!(M)$);
  }
  \fill (C) circle (0.053) (F) circle (0.053);
  \node[fill=white,inner sep=1pt] at (0,0.267) {$C$};
  \draw[line width=0.3pt,black!50] (\R/2,0.049) -- (\R/2,0.588);
  \node[fill=white,inner sep=1pt] at (\R/2,0.776) {$F$};
  \node[anchor=north west,inner sep=1pt] at (\R+0.14,-0.07) {$V$};
  \draw[anteparo] (1.891,-1.715) rectangle (2.029,-0.117);
  \node[anchor=base,text=black!80] at (1.99,-2.101) {anteparo};
  \pic at (A) {brilho};
  \fill (A') circle (0.059);
  \node[above left=0pt] at (A) {$A$};
  \node[above left=-1pt] at (A') {$A'$};
  \pic[xscale=-1,scale=1.022] at (-0.078,-1.009) {olho};
\end{tikzpicture}
